\clearpage
\begin{proofbox}
\textbf{证明一（优先）：距离平方，三步配方。}
取 $A(-a,0),B(a,0)$，$a>0$，$\lambda>0,\lambda\ne1$。设 $P(x,y)$。
\[
\begin{aligned}
 PA=\lambda PB
 &\Longleftrightarrow (x+a)^2+y^2
       =\lambda^2\bigl((x-a)^2+y^2\bigr)\\
 &\Longleftrightarrow
 (1-\lambda^2)(x^2+y^2+a^2)+2a(1+\lambda^2)x=0\\
 &\Longleftrightarrow x^2-2hx+y^2+a^2=0,
 \quad h=\frac{a(\lambda^2+1)}{\lambda^2-1}.
\end{aligned}
\]
所以
\[
 (x-h)^2+y^2=h^2-a^2
 =\frac{a^2\bigl((\lambda^2+1)^2-(\lambda^2-1)^2\bigr)}{(\lambda^2-1)^2}
 =\frac{4a^2\lambda^2}{(\lambda^2-1)^2}>0.
\]
这就是圆心 $C(h,0)$、半径 $R=2a\lambda/|\lambda^2-1|$ 的圆。

\textbf{轨迹必须反向检查。}
上述圆上任一点满足平方距离等式。若 $P=B$，等式将变成 $4a^2=0$，矛盾；
若 $P=A$，将变成 $0=4a^2\lambda^2$，也矛盾。
故 $PA,PB>0$。取非负平方根得到 $PA=\lambda PB$，再除以 $PB$ 即得原比值。
于是\textbf{条件与圆方程互相推出}，轨迹既不缺点，也不多点。
另外 $h^2-R^2=a^2$，故 $|h|>a$，圆心确在线段 $AB$ 外。

\textbf{数值短证。} 当 $a=1,\lambda=2$ 时，直接算
\[
 (x+1)^2+y^2=4\bigl((x-1)^2+y^2\bigr)
 \Longleftrightarrow
 (x-5/3)^2+y^2=16/9.
\]
数值题可以直接配方，不必先背一般公式。

\tcbbreak
\textbf{证明二（几何角度）：内外角平分线给出直径。}
先在直线 $AB$ 上设点 $(x,0)$。条件 $|x+a|=\lambda|x-a|$ 等价于
\[
 x+a=-\lambda(x-a)\quad\text{或}\quad x+a=\lambda(x-a).
\]
分别解这两个一次方程，得到内分点与外分点的横坐标
\[
 x_D=\frac{a(\lambda-1)}{\lambda+1},\qquad
 x_E=\frac{a(\lambda+1)}{\lambda-1}.
\]
$|x_D|<a$，所以 $D$ 在内部；$\lambda>1$ 时 $E$ 在 $B$ 外侧，
$0<\lambda<1$ 时 $E$ 在 $A$ 外侧。两点均满足 $DA/DB=EA/EB=\lambda$。
\begin{center}\begin{tikzpicture}[font=\small]
\draw[black!30,rounded corners=2pt] (0,0) rectangle (14.4,8.2);
\node[anchor=west,font=\small\bfseries] at (.5,7.72) {几何构造：内外角平分线垂直，所以 $P$ 在直径 $DE$ 的圆上};
\begin{scope}[shift={(4.5,2.6)},scale=1.4]
\coordinate (A) at (-1.0000000000,0.0000000000);
\coordinate (B) at (1.0000000000,0.0000000000);
\coordinate (C) at (1.6666666667,0.0000000000);
\coordinate (D) at (0.3333333333,0.0000000000);
\coordinate (E) at (3.0000000000,0.0000000000);
\coordinate (P) at (1.6666666667,1.3333333333);
\coordinate (T) at (3.2666666667,2.1333333333);
\draw[black!30] (C) circle (1.3333333333);\draw[black!55] (-1.5,0)--(3.4,0);
\draw[blue,thick] (A)--(P);\draw[teal!80!black,thick] (B)--(P);\draw[blue,dashed] (P)--(T);\draw[orange!85!black,thick] (D)--(P);\draw[purple,thick,dashed] (P)--(E);\pic[draw,angle radius=3mm] {right angle=D--P--E};\fill[black] (A) circle (1.5pt) node[below=6pt] {$A$};
\fill[black] (B) circle (1.5pt) node[below=6pt] {$B$};
\fill[black] (D) circle (1.5pt) node[below left=6pt] {$D$};
\fill[black] (E) circle (1.5pt) node[below right=6pt] {$E$};
\fill[black] (P) circle (1.5pt) node[above left=10pt] {$P$};
\draw[orange!85!black] (1.2820629745,1.1410314873) arc[start angle=-153.4349488229,delta angle=18.4349488229,radius=0.43];
\draw[orange!85!black] (1.3626107508,1.0292774174) arc[start angle=-135.0000000000,delta angle=18.4349488229,radius=0.43];
\draw[purple] (1.4206991891,0.8413983783) arc[start angle=-116.5650511771,delta angle=71.5650511771,radius=0.55];
\draw[purple] (2.0555753963,0.9444246037) arc[start angle=-45.0000000000,delta angle=71.5650511756,radius=0.55];

\end{scope}
\node[anchor=west] at (.8,6.7) {$PD$ 平分内角，$PE$ 平分外角};\node[anchor=west] at (.8,.6) {$\angle DPE=90^\circ$，故 $PD\perp PE$};
\end{tikzpicture}\end{center}

若 $P$ 满足原比值且不在直线 $AB$ 上，在 $\triangle PAB$ 中，
内角平分线定理说明其内角平分线过 $D$，外角平分线定理说明其外角平分线所在直线过 $E$。
这是因为在对应的内部或外部位置上，给定正比例的分点各自唯一。
内角与邻补外角之和为 $180^\circ$，各取一半，故 $PD\perp PE$。
于是 $\angle DPE=90^\circ$，$P$ 在直径为 $DE$ 的圆上。
共线时直接解距离比，只得到 $D,E$，它们也属于该圆。

\textbf{这条几何路线解释“为什么是这条直径”。}
完整反向包含关系可由证明一回验：
\[
 \frac{x_D+x_E}{2}=h,\qquad
 \frac{|x_E-x_D|}{2}=\frac{2a\lambda}{|\lambda^2-1|}=R.
\]
因此直径圆就是证明一已验证的整个轨迹圆，而不只是一个包含轨迹的圆。

\tcbbreak
\textbf{证明三（已学向量时）：以圆心为原点消去一次项。}
对一般定点，用位置向量 $\mathbf a,\mathbf b,\mathbf c$ 表示 $A,B,C$，令
\[
 \mathbf c=\frac{\mathbf a-\lambda^2\mathbf b}{1-\lambda^2}.
\]
则 $\overrightarrow{CA}=\lambda^2\overrightarrow{CB}$。记 $\mathbf u=\overrightarrow{CP}$，
用 $|\mathbf u-\mathbf v|^2=|\mathbf u|^2-2\mathbf u\cdot\mathbf v+|\mathbf v|^2$ 展开：
\[
\begin{aligned}
 PA^2-\lambda^2PB^2
 &=(1-\lambda^2)PC^2
 -2\mathbf u\cdot\bigl(\overrightarrow{CA}-\lambda^2\overrightarrow{CB}\bigr)
 +CA^2-\lambda^2CB^2\\
 &=(1-\lambda^2)PC^2+CA^2-\lambda^2CB^2.
\end{aligned}
\]
由上述向量公式，$CB=L/|1-\lambda^2|$、$CA=\lambda^2L/|1-\lambda^2|$，故
\[
 CA^2-\lambda^2CB^2=-\frac{\lambda^2L^2}{1-\lambda^2}.
\]
因此
\[
 PA^2-\lambda^2PB^2
 =(1-\lambda^2)\left(PC^2-\frac{\lambda^2L^2}{(1-\lambda^2)^2}\right).
\]
因 $1-\lambda^2\ne0$，左端为零当且仅当 $PC=\lambda L/|1-\lambda^2|$。
$A\ne B$ 保证半径为正；与证明一相同，$P=A,B$ 均不满足平方等式，取正平方根可回到原比值。
按坐标读出 $\mathbf c$ 即得一般坐标圆心公式。

\textbf{边界不除零。} 若 $\lambda=1$，直接用 $PA^2=PB^2$，得到到两定点等距的中垂线。
\end{proofbox}
